\section{Experimental Setup}
\label{sec:experiments}

\subsection{Research questions}
State the questions before presenting results. For example:
\begin{description}[leftmargin=2.4em,style=nextline]
  \item[RQ1:] \placeholder{方法是否优于强基线？}
  \item[RQ2:] \placeholder{每个模块分别带来多少收益？}
  \item[RQ3:] \placeholder{方法在哪些数据、规模或条件下失效？}
\end{description}

\subsection{Datasets and preprocessing}
\placeholder{介绍数据来源、规模、划分、许可、预处理和数据泄漏控制。}

\subsection{Baselines and evaluation metrics}
\placeholder{说明基线选择、公平比较方式、评价指标定义及统计检验。}

\subsection{Implementation details}
\placeholder{报告硬件、软件版本、训练轮数、优化器、学习率、随机种子和调参范围。}

\begin{table}[t]
  \centering
  \caption{Template for dataset or experimental configuration statistics.}
  \label{tab:dataset-statistics}
  \begin{tabular}{lrrr}
    \toprule
    Split & Instances & Features & Labels \\
    \midrule
    Train & \placeholder{--} & \placeholder{--} & \placeholder{--} \\
    Validation & \placeholder{--} & \placeholder{--} & \placeholder{--} \\
    Test & \placeholder{--} & \placeholder{--} & \placeholder{--} \\
    \bottomrule
  \end{tabular}
\end{table}
